
\subsubsection{6.1 Mechanism Validation}

The hypothesis posits a four-step causal chain: (1) task constraints impose architectural invariants, (2) NFN encoders preserve task-relevant features, (3) hierarchical pooling exposes global structure, (4) Transformer potentially learns cross-architecture relational structure.

\textbf{Steps 1-3 validated on synthetic data:} CKA same-task 0.82 (Step 1), reconstruction accuracy 68\% (Step 2), WCSS clustering validated with large effect size (Step 3).

\textbf{Step 4 inferred but unconfirmed:} Clustering success suggests Transformer may contribute to alignment, but no architecture token ablation was performed. Cannot claim ''Transformer discovers cross-architecture correspondences'' without ablation evidence.

\subsubsection{6.2 Unexpected Findings}

Three findings on synthetic data exceeded expectations: (1) effect size 1.45 vs planned 0.5, (2) CKA 0.82 vs threshold 0.6, (3) reconstruction accuracy marginal 68\% vs target 70\%.

\textbf{Competing explanations for (1) and (2):} Task structure stronger than expected (supports hypothesis) vs synthetic data artifact (threatens validity). Evidence supports artifact hypothesis: synthetic models may embed task signals more cleanly than real checkpoints. Real ModelZooDataset validation required.

\textbf{Explanation for (3):} Early stopping artifact (10 epochs vs 200 planned). Reconstruction loss still descending. Full training expected to improve accuracy to 70-75\%.

\subsubsection{6.3 Limitations}

\textbf{L1: Synthetic Dataset (CRITICAL):} All results from synthetic data, not real Zenodo ModelZooDataset downloads. External validity unconfirmed. CKA scores and effect sizes may be inflated. Real dataset validation required before publication.

\textbf{L4: Reduced Training Scale:} 10 epochs (PoC) vs 200 epochs (planned). Reconstruction accuracy marginal (68\% vs 70\% target). Full-scale training recommended.

\textbf{L5: Architecture Token Ablation Deferred:} No ablation study removing architecture-type token embeddings. Cannot quantify Transformer contribution to clustering. Mechanism claim weakened.

\textbf{L7: Generative Models Excluded:} GANs/Diffusion models excluded by design. Scope limited to discriminative architectures.

\textbf{L8: Fine-Grained Editing Operations Excluded:} Lossy pooling discards neuron-level details required for task arithmetic and model merging. Method suitable for inference only.

\subsubsection{6.4 Future Work}

\textbf{Priority 1 (Critical for Publication):} Real dataset validation. Download full ModelZooDataset from Zenodo, re-run CKA gate + WCSS test on real checkpoints. Acceptance criteria: CKA same-task $> 0.6$ AND WCSS Cohen's $d > 0.5$ on real data. Timeline: 4 days. Resource: $1\times$ V100 GPU (\$50).

\textbf{Priority 2 (Marginal for Reconstruction):} Full-scale training. Train hierarchical VAE for 200 epochs on $2\times$ V100 GPUs to improve reconstruction accuracy from 68\% to 70-75\%. Timeline: 7 days. Resource: $2\times$ V100 (\$700).

\textbf{Priority 3 (Mechanism Refinement):} Architecture token ablation. Retrain Transformer without architecture-type tokens, measure clustering degradation. Acceptance: If degradation $\geq 15pp$, tokens critical. If $< 5\%$, simplify to pooling-only. Timeline: 2 days. Resource: $1\times$ V100 (\$50).

\textbf{Extension 1:} Baseline comparison. Compare hierarchical VAE vs UNF encoders + architecture-conditioned MLP and SANE to demonstrate performance improvement.

\textbf{Extension 2:} MLP/ViT encoder implementation. Expand coverage from 2 to 4 architecture families (285 MLPs, 225 ViTs in dataset).

\textbf{Extension 3:} Zero-shot task prediction on Hugging Face. Collect 100 models with corrupted metadata, test zero-shot task prediction.

